\documentclass[11pt]{article}

\usepackage[margin=1in]{geometry}
\usepackage{amsmath}
\usepackage{amssymb}
\usepackage[hidelinks]{hyperref}

\title{Memory-Efficient MCWF Squeezing Analysis}
\author{}
\date{}

\begin{document}
\maketitle

\section{Conclusion}

The current squeezing analysis does not require cross-trajectory products.
Its only ensemble-wide dependencies are the averaged $J$ direction used to
define the $S$ basis and the averaged $J$ and $S$ directions used to define
the squeezing operators.  Consequently, the calculation can be performed
with one input file in memory at a time using three streaming passes.  A more
fundamental redesign can eliminate stored trajectories entirely by
accumulating complete basis-independent first- and second-moment tensors while
each trajectory is simulated.

\section{Three Streaming Passes}

An ensemble average only requires a running sum,
\begin{equation}
  \overline{X}(t_k)
  = \frac{1}{R}\sum_{r=1}^{R}X^{(r)}(t_k),
\end{equation}
not simultaneous access to all trajectories.  The current basis dependencies
can therefore be respected as follows.

\begin{enumerate}
  \item Load one ensemble file, compute the raw $J$ moments of its
  trajectories, add them to running sums, and release the file.  Repeat for
  every file, then construct the ensemble-averaged $J$ direction and angles.

  \item Reload one file at a time, compute raw $S$ moments using the now-fixed
  ensemble $J$ angles, accumulate them, and release the file.  After all files
  have been processed, construct the ensemble-averaged $S$ direction and
  angles.

  \item Reload one file at a time, compute the raw squeezing moments using the
  fixed ensemble $J$ and $S$ angles, accumulate them, and release the file.
  Construct the covariance matrix and squeezing only after the final average
  has been formed.
\end{enumerate}

With the present pickle format, one complete array-job file must still be
loaded at once.  The peak ensemble memory is nevertheless reduced from all
four files to one file.  The moment-extraction functions would need streaming
accumulators rather than their current interface, which collects all
per-trajectory results before averaging.

\section{Basis-Independent Sufficient Moments}

The three-pass method can be avoided by accumulating a complete set of fixed
basis moments during simulation.  For each physical group $g$, define the nine
collective matrix-unit operators
\begin{equation}
  A_{\mu\nu,g}
  = \sum_{i\in g} |\mu_i\rangle\langle\nu_i|,
  \qquad
  \mu,\nu\in\{u,d,e\}.
\end{equation}
Let a single index $p$ enumerate the nine ordered pairs $(\mu,\nu)$.  For each
trajectory $r$ and saved time $t_k$, compute the individually normalized
moments
\begin{align}
  m_{p,g}^{(r)}(t_k)
  &=
  \frac{
    \langle\psi_r(t_k)|A_{p,g}|\psi_r(t_k)\rangle
  }{Z_r(t_k)},
  \\
  G_{pq,g}^{(r)}(t_k)
  &=
  \frac{
    \langle A_{p,g}\psi_r(t_k)
    |A_{q,g}\psi_r(t_k)\rangle
  }{Z_r(t_k)}.
\end{align}
Accumulate only
\begin{equation}
  \overline{m}_{p,g}
  = \frac{1}{R}\sum_r m_{p,g}^{(r)},
  \qquad
  \overline{G}_{pq,g}
  = \frac{1}{R}\sum_r G_{pq,g}^{(r)},
\end{equation}
and discard the trajectory sector blocks.  The matrix
$\overline{G}_{pq,g}$ must remain complex; storing only its real part would
not be sufficient for later basis rotations.

These accumulators can be produced either after each trajectory finishes or
directly at each saved snapshot.  Computing them at snapshot time avoids even
retaining all snapshots of one trajectory.  With multiprocessing, a worker
can return one accumulated tensor, or maintain a worker-local sum and return
it when its assigned trajectories are complete.

\section{Reconstruction of the Current Analysis}

The ensemble $J$ moments are linear combinations of
$\overline{m}$.  For example, within one group,
\begin{align}
  J_x &= \frac{\overline{m}_{ed}+\overline{m}_{de}}{2},
  &
  J_y &= \frac{\overline{m}_{ed}-\overline{m}_{de}}{2i},
  \\
  J_z &= \frac{\overline{m}_{ee}-\overline{m}_{dd}}{2},
  &
  N_e &= \overline{m}_{ee},
  \qquad
  N_J = \overline{m}_{dd}+\overline{m}_{ee}.
\end{align}
The normalized $J$ direction and the angles $\theta_J,\phi_J$ are then
constructed exactly as in the current moment pipeline.

Once the ensemble $J$ angles are known, the effective $S$ basis and its three
collective spin operators are fixed.  Every such operator is a linear
combination of the same $A_{p,g}$, so its ensemble expectation value follows
directly from $\overline{m}$.  These values determine the ensemble $S$
direction and angles $\theta_S,\phi_S$ without revisiting any trajectory.

The dressed states $|c_g\rangle$, $|j_g\rangle$, and $|s_g\rangle$ can then be
constructed from the ensemble $J$ and $S$ angles.  Each of the four current
fluctuation operators has an expansion
\begin{equation}
  O_{a,g}=\sum_p c_{ap,g}A_{p,g}.
\end{equation}
The required squeezing moments are recovered as
\begin{align}
  \mu_{a,g}
  &= \sum_p c_{ap,g}\,\overline{m}_{p,g},
  \\
  M_{ab,g}
  &= \operatorname{Re}
  \sum_{p,q}
  c_{ap,g}^{*}\,\overline{G}_{pq,g}\,c_{bq,g}.
\end{align}
The mean population $N_{c,g}$ is another linear combination of
$\overline{m}_{p,g}$ obtained from the projector
$|c_g\rangle\langle c_g|$.  The existing solver-independent squeezing code
can then construct
\begin{equation}
  C_{ab,g}=M_{ab,g}-\mu_{a,g}\mu_{b,g},
\end{equation}
its minimum eigenvalue, and $\xi_g^2$ without modification to the physical
definition.

\section{Why the Reconstruction Is Exact}

At a fixed saved time, the coefficients $c_{ap,g}$ are obtained from the
ensemble-averaged angles and are therefore identical for every trajectory.
Linearity gives
\begin{equation}
  \frac{1}{R}\sum_r
  c_a^\dagger G_g^{(r)}c_b
  =
  c_a^\dagger
  \left(\frac{1}{R}\sum_r G_g^{(r)}\right)c_b
  =
  c_a^\dagger\overline{G}_g c_b.
\end{equation}
The same identity holds for every first moment.  Thus rotating the averaged
fixed-basis tensors after the ensemble angles are known gives the same result
as rotating and evaluating every stored trajectory afterward, up to ordinary
floating-point summation order.

Exact equivalence requires the following invariants:
\begin{itemize}
  \item normalize each trajectory by its own $Z_r(t_k)$ before accumulation;
  \item retain the complete complex $9\times9$ matrix
  $\overline{G}_{pq,g}$;
  \item derive $J$ and $S$ angles from ensemble-averaged first moments;
  \item construct covariance matrices and squeezing only after ensemble
  averaging, rather than averaging per-trajectory squeezing values.
\end{itemize}

The current analysis is group-local, so one $9\times9$ tensor per group is
sufficient.  A future cross-group squeezing observable would additionally
require cross-group blocks
$\overline{G}_{(p,g)(q,h)}$.  Diagnostics such as the jump rate can be
accumulated directly during simulation or reconstructed from the corresponding
fixed-basis correlations.

For two groups and 500 saved times, the fixed-basis averages contain only
$2\times500\times9$ complex first moments and
$2\times500\times9\times9$ complex second moments, approximately
$1.4\,\mathrm{MiB}$ of numerical data.  Their size is independent of the
number of trajectories, replacing multi-gigabyte temporary trajectory files
with a small set of ensemble sufficient statistics.

\end{document}
